\documentclass[12pt]{article}
\usepackage[T1]{fontenc}
\usepackage[utf8]{inputenc}
\usepackage{lmodern}
\usepackage{textcomp}
\usepackage{enumitem}
\usepackage{geometry}
\usepackage{graphicx}
\usepackage{amsmath, amssymb}
\geometry{margin=1in}
\begin{document}
\begin{center}
History - Undergraduate Level Assessment 9 \
Total Marks: 40 \
Answer all questions.
\end{center}

\section*{Multiple Choice (10 marks)}
\begin{enumerate}[label=\arabic*.]
\item The most prominent feature of the food resource base of post-Pleistocene Europe was its:
\begin{enumerate}[label=(\alph*)]
    \item dependence on megafauna.
    \item hunting and gathering.
    \item dependence on fur-bearing animals.
    \item diversity.
\end{enumerate}
\item The fossil evidence indicates that by 1.8 million years ago, a new hominid known as \_\_\_\_\_\_\_ had evolved.
\begin{enumerate}[label=(\alph*)]
    \item Australopithecus afarensis
    \item Homo naledi
    \item Homo neanderthalensis
    \item Neanderthals
    \item Homo floresiensis
\end{enumerate}
\item What do we know about the people who built Stonehenge based on the analysis of animal teeth from a nearby village?
\begin{enumerate}[label=(\alph*)]
    \item They had no domesticated animals but hunted wild animals for meat.
    \item They had domesticated animals, but they did not consume their meat.
    \item They bred and raised a specific species of deer for food.
    \item They were vegetarian and did not consume meat.
    \item They brought animals to the site from outside of the region.
\end{enumerate}
\item What occupation date do most archaeologists agree on for Australia?
\begin{enumerate}[label=(\alph*)]
    \item soon after 60,000 years ago
    \item approximately 50,000 years ago
    \item no earlier than 100,000 years ago
    \item soon after 30,000 years ago
    \item no later than 70,000 years ago
\end{enumerate}
\item Taphonomy is:
\begin{enumerate}[label=(\alph*)]
    \item the study of plant and animal life from a particular period in history.
    \item a study of the spatial relationships among artifacts, ecofacts, and features.
    \item the study of human behavior and culture.
    \item a technique to determine the age of artifacts based on their position in the soil.
    \item a dating technique based on the decay of a radioactive isotope of bone.
\end{enumerate}
\end{enumerate}

\section*{True / False (10 marks)}
\textit{Answer True or False}
\begin{enumerate}[label=\arabic*.]
\setcounter{enumi}{5}
\item True or False: He surveyed editions of medieval French texts that were produced with the stemmatic method, and found that textual critics tended overwhelmingly to produce trees divided into just two branches. He concluded that this outcome was unlikely to have occurred by chance, and that therefore, the method was tending to produce bipartite stemmas regardless of the actual history of the witnesses. He suspected that editors tended to favor trees with two branches, as this would maximize the opportunities for editorial judgment (as there would be no third branch to "break the tie" whenever the witnesses disagreed).
\item True or False: To address the problem of slowed timekeeping in the pressure head of the inflow water clock, Zhang was the first in China to install an additional tank between the reservoir and inflow vessel. Zhang also invented a seismometer (Houfeng didong yi 候风地动仪) in a different answer to detect the exact cardinal or ordinal direction of earthquakes from hundreds of kilometers away.
\item True or False: The fourth a cappella musical to appear Off-Broadway, In Transit, premiered 5 October 2010 and was produced by Primary Stages with book, music, and lyrics by Kristen Anderson-Lopez, James-Allen Ford, Russ Kaplan, and Sara Wordsworth.
\item True or False: Feminist anthropology is a four field approach to anthropology (archeological, biological, cultural, linguistic) that seeks to reduce male bias in research findings, anthropological hiring practices, and the scholarly production of knowledge.
\item True or False: Twelve days after Chinese troops entered the war, Stalin allowed the Soviet Air Force to provide air cover, and supported more aid to China. After decimating the ROK II Corps at the Battle of Onjong, the first confrontation between Chinese and U.S. military occurred on 1 November 1950; deep in North Korea, thousands of soldiers from the PVA 39 th Army encircled and attacked the U.
\end{enumerate}

\section*{Long Form Response (20 marks)}
\textit{Answer each question with a clear and complete explanation.}
\begin{enumerate}[label=\arabic*.]
\setcounter{enumi}{10}
\item Read the following passage and answer the question: A "lock-in" is when a pub owner lets drinkers stay in the pub after the legal closing time, on the theory that once the doors are locked, it becomes a private party rather than a pub. Patrons may put money behind the bar before official closing time, and redeem their drinks during the lock-in so no drinks are technically sold after closing time. The origin of the British lock-in was a reaction to 1915 changes in the licensing laws in England and Wales, which curtailed opening hours to stop factory workers from turning up drunk and harming the war effort. Since 1915, the UK licensing laws had changed very little, with comparatively early closing times. The tradition of the lock-in therefore remained. Since the implementation of Licensing Act 2003, premises in England and Wales may apply to ... Question: What action by a pub owner can result in his prosecution?
\item Read the following passage and answer the question: The history of the area that now constitutes Himachal Pradesh dates back to the time when the Indus valley civilisation flourished between 2250 and 1750 BCE. Tribes such as the Koilis, Halis, Dagis, Dhaugris, Dasa, Khasas, Kinnars, and Kirats inhabited the region from the prehistoric era. During the Vedic period, several small republics known as "Janapada" existed which were later conquered by the Gupta Empire. After a brief period of supremacy by King Harshavardhana, the region was once again divided into several local powers headed by chieftains, including some Rajput principalities. These kingdoms enjoyed a large degree of independence and were invaded by Delhi Sultanate a number of times. Mahmud Ghaznavi conquered Kangra at the beginning of the 10 th century. Timur and Sikander Lodi al... Question: What tribes inhibited the area that now constitutes Himachal Pradesh?
\end{enumerate}

\vfill
\noindent\textit{End of Paper}
\end{document}